\documentclass{article}
\usepackage{neurips_2024}

\usepackage[utf8]{inputenc}
\usepackage[T1]{fontenc}
\usepackage{hyperref}
\usepackage{url}
\usepackage{booktabs}
\usepackage{amsfonts}
\usepackage{amsmath}
\usepackage{amssymb}
\usepackage{amsthm}
\usepackage{nicefrac}
\usepackage{microtype}
\usepackage{xcolor}
\usepackage{algorithm}
\usepackage{algorithmic}
\usepackage{enumitem}
\usepackage{tikz}
\usepackage{stmaryrd}       % for \llbracket, \rrbracket

% Theorem environments
\newtheorem{definition}{Definition}
\newtheorem{assumption}{Assumption}
\newtheorem{proposition}{Proposition}
\newtheorem{theorem}{Theorem}
\newtheorem{lemma}{Lemma}
\newtheorem{corollary}{Corollary}
\newtheorem{remark}{Remark}
\newtheorem{example}{Example}

% Operator shortcuts
\DeclareMathOperator*{\argmin}{arg\,min}
\DeclareMathOperator*{\argmax}{arg\,max}
\DeclareMathOperator{\pre}{pre}
\DeclareMathOperator{\eff}{eff}
\DeclareMathOperator{\dom}{dom}
\DeclareMathOperator{\img}{img}
\newcommand{\calS}{\mathcal{S}}
\newcommand{\calA}{\mathcal{A}}
\newcommand{\calO}{\mathcal{O}}
\newcommand{\calC}{\mathcal{C}}
\newcommand{\calD}{\mathcal{D}}
\newcommand{\calE}{\mathcal{E}}
\newcommand{\calF}{\mathcal{F}}
\newcommand{\calG}{\mathcal{G}}
\newcommand{\calH}{\mathcal{H}}
\newcommand{\calI}{\mathcal{I}}
\newcommand{\calK}{\mathcal{K}}
\newcommand{\calL}{\mathcal{L}}
\newcommand{\calM}{\mathcal{M}}
\newcommand{\calN}{\mathcal{N}}
\newcommand{\calP}{\mathcal{P}}
\newcommand{\calR}{\mathcal{R}}
\newcommand{\calT}{\mathcal{T}}
\newcommand{\calU}{\mathcal{U}}
\newcommand{\calV}{\mathcal{V}}
\newcommand{\calW}{\mathcal{W}}
\newcommand{\calX}{\mathcal{X}}
\newcommand{\bbR}{\mathbb{R}}
\newcommand{\bbN}{\mathbb{N}}
\newcommand{\bbZ}{\mathbb{Z}}
\newcommand{\ind}{\mathbb{1}}
\newcommand{\Obj}{\mathsf{Obj}}
\newcommand{\Pred}{\mathsf{Pred}}
\newcommand{\strips}{\textsc{Strips}}
\newcommand{\pddl}{\textsc{Pddl}}

\title{Object-Centric World Model Synthesis via Factored Program Induction}

\author{
  David \\
  \texttt{TODO} \\
}

\begin{document}

\maketitle

\begin{abstract}
We extend WorldCoder~\cite{tang2024worldcoder}, a system that synthesizes
world models as Python programs via counterexample-guided inductive synthesis
(CEGIS), in four directions.  First, we replace its monolithic transition
function with an \emph{object-centric OOP factorization}: each object type is
a class whose \texttt{respond()} method encodes how instances of that type
react to actions in context.  This reduces the synthesis search space from
monolithic to factored, isolates regressions per type, and enables cross-domain
transfer via class inheritance.  Second, we replace WorldCoder's binary
known/unknown partition with a \emph{continuous epistemic model} that tracks
observation count, context diversity, and effect consistency per
(object-type, action) cell, driving exploration that actively seeks diverse
contexts rather than merely covering the action space.  Third, we make
transition functions \emph{object-centric}: instead of a single
$\hat{T}(s,a) \to s'$, each object type owns its dynamics and receives
relational context (adjacency, selection, co-location) as input, enabling the
system to learn interaction effects that depend on spatial relationships between
objects.  Fourth, we replace WorldCoder's stateless backprompting with an
\emph{agentic synthesis backend}: the LLM operates as an autonomous agent with
persistent workspace access, incremental file editing, and autonomous test
execution, implementing the CEGIS loop through interactive tool use rather
than prompt engineering.  We evaluate on WorldCoder's benchmarks (Sokoban,
MiniGrid, AlfWorld) and on ARC-AGI-3, a pixel-observation gridworld benchmark
with anonymous actions.  Experiments demonstrate that OOP factoring reduces
synthesis token cost, the continuous epistemic model improves sample efficiency,
agentic synthesis reaches higher accuracy than stateless backprompting, and
learned object classes transfer across domain variants.
\end{abstract}


% ============================================================================
% SECTION 1: INTRODUCTION
% ============================================================================

\section{Introduction}
\label{sec:intro}

Model-based reinforcement learning (MBRL) builds an internal model of the
environment's dynamics and uses it for planning.  Most MBRL systems learn
world models as neural networks trained via gradient descent on prediction
losses~\cite{ha2018world, hafner2020dreamer, schrittwieser2020muzero}.  These
models are sample-hungry, opaque, and domain-specific: a world model trained on
one Atari game does not transfer to another.

WorldCoder~\cite{tang2024worldcoder} introduced a radically different approach:
use a large language model (LLM) to \emph{synthesize} the world model as an
executable Python program.  The program takes as input a structured state
(a set of typed objects with attributes) and an action, and returns the
predicted next state.  Synthesis follows a counterexample-guided inductive
synthesis (CEGIS) loop: the LLM generates a candidate program, the system
verifies it against observed transitions, and counterexamples (transitions
where the program's prediction disagrees with reality) are fed back to the
LLM for refinement.  An optimism constraint ($\phi_2$) drives exploration by
requiring that the current model admit a plan leading to reward.

WorldCoder achieves remarkable sample efficiency---roughly 50 environment
interactions to learn accurate world models for Sokoban, versus $>$1M steps
for PPO or DreamerV3---and demonstrates transfer across MiniGrid task variants.
However, four limitations constrain its scalability:

\begin{enumerate}[nosep]
    \item \textbf{Monolithic synthesis.}  The LLM synthesizes a single
          \texttt{transition\_function(state, action)} that handles all object
          types and interactions in one function.  Fixing one object's behavior
          risks breaking others.  The search space grows combinatorially with
          the number of object types and actions.

    \item \textbf{Binary exploration model.}  WorldCoder partitions knowledge
          into known (verified correct) and unknown (not yet tested).  This
          conflates ``never tested'' with ``tested once in one context,''
          missing the distinction between an effect that was absent (true
          non-effect) and an effect that was blocked by context (false negative).

    \item \textbf{No structural transfer.}  Each domain starts synthesis from
          scratch.  A \texttt{Wall} class learned in Sokoban cannot be reused in
          a new domain---the entire transition function must be re-synthesized.

    \item \textbf{Stateless synthesis.}  Each CEGIS iteration is a stateless
          LLM API call: the model receives the full prompt (prior code plus
          counterexamples) and returns a complete candidate program.  The LLM
          has no persistent memory across iterations, cannot incrementally edit
          code, and cannot autonomously run tests or inspect intermediate
          results.  This limits the LLM to shallow, single-pass reasoning
          per iteration.
\end{enumerate}

We address these limitations with four contributions:

\begin{enumerate}[nosep]
    \item \textbf{OOP-factored synthesis} (\S\ref{sec:oop_state}--\S\ref{sec:oop_transition}):
          The world model is a set of typed object classes.  Each class has a
          \texttt{respond(action, context)} method encoding how that type reacts
          to actions.  The CEGIS loop synthesizes one class at a time, and
          counterexamples are attributed to specific types.  We prove that this
          factorization reduces the CEGIS sample complexity from monolithic to
          additive (Theorem~\ref{thm:factored_bound}).

    \item \textbf{Continuous epistemic model} (\S\ref{sec:epistemic}):
          Each (object-type, action) cell tracks observation count, context
          diversity, effect consistency, and conditional effect flags.
          Exploration priority is a continuous score rather than a binary
          known/unknown label, driving the agent to seek diverse contexts and
          actively test hypotheses about conditional preconditions.

    \item \textbf{Object-centric transitions with relational context}
          (\S\ref{sec:oop_transition}):
          Each object's \texttt{respond()} method receives a context object
          providing relational queries (adjacency, co-location, selection state).
          This enables the system to learn interaction effects that depend on
          spatial relationships between objects, without hard-coding the
          relational predicates in advance.

    \item \textbf{Agentic synthesis backend} (\S\ref{sec:agentic}):
          The CEGIS loop is implemented by an autonomous coding agent with
          persistent workspace access, incremental file editing, and
          interactive test execution.  The agent maintains full conversation
          history across edit-test-debug cycles within a single synthesis
          invocation, enabling multi-step reasoning, targeted fixes, and
          strategic debugging that stateless backprompting cannot achieve.
\end{enumerate}

We evaluate on WorldCoder's three benchmarks (Sokoban, MiniGrid, AlfWorld)
plus ARC-AGI-3~\cite{arcagi3}, a pixel-observation gridworld benchmark where
the agent receives raw 64$\times$64 RGB frames and has no knowledge of action
semantics.


% ============================================================================
% SECTION 2: PRELIMINARIES
% ============================================================================

\section{Preliminaries}
\label{sec:preliminaries}

\subsection{Object-Oriented MDP}
\label{sec:oo_mdp}

Following WorldCoder~\cite{tang2024worldcoder} and the OO-MDP formalism
of Diuk et al.~\cite{diuk2008oomdp}, we model the environment as:

\begin{definition}[Object-Oriented MDP]
\label{def:oomdp}
An \textbf{object-oriented MDP} is a tuple
$\calM = (\calC, \calA, T, R, s_0)$, where:
\begin{itemize}[nosep]
    \item $\calC = \{C_1, \ldots, C_K\}$ is a finite set of \textbf{object
          classes}.  Each class $C_k$ is defined by a set of \textbf{attributes}
          $\text{Att}(C_k)$.
    \item Each state $s = \{o_1, \ldots, o_n\}$ is a set of \textbf{object
          instances}, where each $o_i$ is an instance of some class $C_{k_i}$
          with attribute values $o_i.\text{att}$ for each $\text{att} \in
          \text{Att}(C_{k_i})$.
    \item $\calA = \{a_1, \ldots, a_N\}$ is a finite action set.
    \item $T : \calS \times \calA \to \calS$ is the deterministic transition
          function.
    \item $R : \calS \times \calA \times \calS \to \bbR$ is the reward function.
    \item $s_0 \in \calS$ is the initial state.
\end{itemize}
\end{definition}

WorldCoder represents states as object-oriented MDPs but synthesizes
transitions as a monolithic function.  Our key departure is distributing
transition logic to the objects themselves.

\subsection{WorldCoder: CEGIS for World Models}
\label{sec:worldcoder}

WorldCoder synthesizes two Python subroutines:
\begin{itemize}[nosep]
    \item $\texttt{transition\_function}(s, a) \to s'$: predicts the next state.
    \item $\texttt{reward\_function}(\text{context})(s, a, s') \to (r, \text{done})$:
          predicts reward and termination.
\end{itemize}
Synthesis follows a CEGIS loop with two constraints:

\begin{definition}[WorldCoder Constraints]
\label{def:wc_constraints}
Given replay buffer $\calD_t = \{(s_i, a_i, s'_i)\}_{i=1}^t$ and candidate
program $\hat{T}$:
\begin{align}
    \phi_1(\calD_t, \hat{T}) &= \forall (s, a, s') \in \calD_t :
        \hat{T}(s, a) = s'
        \tag{data consistency} \\
    \phi_2(s_0, \hat{T}) &= \exists a_1, \ldots, a_\ell :
        \text{the trajectory under } \hat{T} \text{ from } s_0
        \text{ reaches a rewarding state}
        \tag{optimism}
\end{align}
\end{definition}

The optimism constraint $\phi_2$ drives exploration: the model must be
consistent with observed data ($\phi_1$) \emph{and} optimistic about
unobserved regions ($\phi_2$), causing the agent to explore when its current
model predicts a path to reward through unknown territory.

\begin{definition}[WorldCoder Sample Complexity {\cite[Theorem~2.4]{tang2024worldcoder}}]
\label{def:wc_bound}
Let $\calW$ be the hypothesis class of candidate programs with logical
dimensionality $K_{\bot\!\!\bot}(\calW)$ (the maximum number of counterexamples
needed to eliminate all incorrect hypotheses), and let $D$ be the environment
diameter.  Then WorldCoder's total interaction count is at most:
\begin{equation}
    D \cdot (K_{\bot\!\!\bot}(\calW) + 1).
\end{equation}
\end{definition}

This bound is tight: $D$ interactions are needed to traverse the state space
between counterexample-producing states, and $K_{\bot\!\!\bot}(\calW)$
counterexamples are needed to converge to the correct program.  Our
contribution is showing that OOP factorization reduces
$K_{\bot\!\!\bot}(\calW)$.


% ============================================================================
% SECTION 3: OOP STATE REPRESENTATION
% ============================================================================

\section{Object-Oriented State Representation}
\label{sec:oop_state}

\subsection{Object Types as Classes}
\label{sec:type_hierarchy}

\begin{definition}[Object Type]
\label{def:object_type}
An \textbf{object type} (class) $\tau$ is a tuple:
\begin{equation}
    \tau = (\text{name}, \; \tau_{\text{parent}}, \;
           \calP_\tau^{\text{atomic}}, \;
           \calP_\tau^{\text{derived}}, \;
           f_\tau^{\text{respond}}),
\end{equation}
where:
\begin{itemize}[nosep]
    \item $\text{name} \in \Sigma^*$ is the class name
          (e.g., \texttt{Wall}, \texttt{Player}, \texttt{Box}).
    \item $\tau_{\text{parent}}$ is the parent type ($\bot$ for the root
          \texttt{GameObject}).
    \item $\calP_\tau^{\text{atomic}}$ is the set of \textbf{atomic predicates}
          (mutable state): position, color, alive, etc.
    \item $\calP_\tau^{\text{derived}}$ is the set of \textbf{derived predicates}
          (computed from atomic predicates, never stored):
          $\mathsf{can\_move}()$, $\mathsf{is\_blocked}()$, etc.
    \item $f_\tau^{\text{respond}} : (\text{self}, \text{action},
          \text{context}) \to \text{self}'$ is the \textbf{response method}:
          given an action and relational context, compute this object's updated
          state.  Subclasses override this method.
\end{itemize}
\end{definition}

\begin{definition}[Type Hierarchy]
\label{def:type_tree}
The \textbf{type hierarchy} $\calH = (\calV, E)$ is a rooted tree where
$\calV$ is the set of object types, $E$ connects each type to its parent, and
the root is \texttt{GameObject}.  A child type inherits all predicates and
methods from its parent and may override any of them.
\end{definition}

\begin{example}[Sokoban hierarchy]
\label{ex:sokoban}
\begin{verbatim}
GameObject (root)
  +-- StaticObject        [respond = no-op]
  |     +-- Wall
  |     +-- Goal
  +-- DynamicObject       [respond = move by delta if unblocked]
        +-- Player        [respond = move in action direction]
        +-- Box           [respond = move when pushed by player]
\end{verbatim}
\texttt{StaticObject.respond()} returns self unchanged.
\texttt{Box.respond()} checks whether the player is adjacent and pushing,
then moves in the push direction if unblocked.
\end{example}

\subsection{Atomic vs.\ Derived Predicates}
\label{sec:predicates}

\begin{definition}[Atomic Predicate]
\label{def:atomic}
An \textbf{atomic predicate} is a mutable property stored explicitly as part
of an object's state.  Only atomic predicates are modified by
$f^{\text{respond}}$.  Examples: $\mathsf{position}(o)$,
$\mathsf{color}(o)$, $\mathsf{alive}(o)$.
\end{definition}

\begin{definition}[Derived Predicate]
\label{def:derived}
A \textbf{derived predicate} is computed deterministically from atomic
predicates.  It is never stored, only recomputed on demand.  Examples:
\begin{align}
    \mathsf{blocked}(o, d) &= \exists o' : \mathsf{position}(o')
        = \mathsf{position}(o) + \Delta(d) \;\wedge\; o'.\text{is\_solid} \\
    \mathsf{adjacent}(o_1, o_2) &= \|\mathsf{position}(o_1) -
        \mathsf{position}(o_2)\|_\infty \leq 1
\end{align}
\end{definition}

\begin{remark}[Why the distinction matters]
After $f^{\text{respond}}$ updates atomic predicates, all derived predicates
are automatically consistent.  This eliminates the frame problem for derived
predicates---analogous to the distinction between fluents and derived predicates
in PDDL~\cite{fox2003pddl2} and between base tables and views in databases.
\end{remark}

\subsection{Structured State}
\label{sec:structured_state}

\begin{definition}[Structured State]
\label{def:structured_state}
The \textbf{structured state} is:
\begin{equation}
    \bar{s} = \{(o_i, \tau_i, \mathbf{p}_i)\}_{i=1}^{M},
\end{equation}
where $\{o_i\}$ are object instances with unique identifiers (track IDs),
$\{\tau_i\}$ are their assigned types, and $\{\mathbf{p}_i\}$ are their
atomic predicate values.
\end{definition}

\begin{definition}[Perception Function]
\label{def:perception}
The \textbf{perception function} $\text{Perceive} : \calS_{\text{raw}} \to
\bar{\calS}$ maps raw observations to structured states.  For symbolic domains
(Sokoban, MiniGrid, AlfWorld), this is a direct read of the environment's
object representation---identical to WorldCoder's setup.  Perception is
\textbf{not} part of the synthesis target; it is a fixed, deterministic
preprocessing step.

For pixel-observation domains (ARC-AGI-3), the environment does not provide
structured state, so we implement a deterministic perception pipeline as a
compatibility layer: connected-component detection $\to$ sprite registry
matching $\to$ cross-frame tracking $\to$ type assignment.  This is not a
contribution of our approach but a necessary adapter for evaluating on
pixel-observation benchmarks.  Details are in Section~\ref{sec:arc_perception}.
\end{definition}

\begin{remark}[Perception must be excluded from synthesis]
\label{rem:perception}
On ARC-AGI-3 games, when the LLM was asked to synthesize a
\texttt{perceive()} function alongside transition logic, it spent all
synthesis iterations debugging perception and achieved 0\% transition
accuracy.  When perception was fixed and the LLM synthesized only
\texttt{respond()} methods, accuracy reached 100\% within 4 iterations.
This motivates our decision to treat perception as a fixed preprocessing
step on all benchmarks.
\end{remark}


% ============================================================================
% SECTION 4: OBJECT-CENTRIC TRANSITION
% ============================================================================

\section{Object-Centric Transition Model}
\label{sec:oop_transition}

\subsection{The \texttt{respond()} Method}
\label{sec:respond}

The central novelty of our approach is that each object type owns its
transition logic via the \texttt{respond()} method.  WorldCoder synthesizes:
\begin{equation}
    \hat{T}_{\text{mono}}(s, a) \to s' \qquad \text{(monolithic)}
\end{equation}
We synthesize, for each object type $\tau_k$:
\begin{equation}
    f_{\tau_k}^{\text{respond}}(o, a, \text{ctx}) \to o' \qquad
    \text{(per-type)}
\end{equation}
where $\text{ctx}$ provides relational queries over the current state.

\begin{definition}[Relational Context]
\label{def:context}
The \textbf{relational context} $\text{ctx}(\bar{s})$ is a read-only view of
the structured state providing:
\begin{itemize}[nosep]
    \item $\text{ctx}.\text{objects\_at}(p) \to \{o_j\}$: objects at position $p$.
    \item $\text{ctx}.\text{adjacent}(o, d) \to o_j \cup \{\bot\}$: object
          adjacent to $o$ in direction $d$, or $\bot$ if empty.
    \item $\text{ctx}.\text{player}() \to o_j$: the player object (if identified).
    \item $\text{ctx}.\text{selected}() \to o_j \cup \{\bot\}$: the currently
          selected object, if any.
    \item $\text{ctx}.\text{of\_type}(\tau) \to \{o_j\}$: all instances of type $\tau$.
\end{itemize}
The context is immutable within a single transition step: objects read the
pre-action state, not partially updated state.
\end{definition}

\begin{definition}[Object-Centric Transition]
\label{def:oop_transition}
Given structured state $\bar{s} = \{(o_i, \tau_i, \mathbf{p}_i)\}$ and
action $a$, the \textbf{object-centric transition} proceeds:
\begin{enumerate}[nosep]
    \item Compute context: $\text{ctx} \leftarrow \text{ctx}(\bar{s})$.
    \item For each object $o_i$ in \textbf{causal order}:
          $\mathbf{p}'_i \leftarrow f_{\tau_i}^{\text{respond}}(o_i, a, \text{ctx})$.
    \item Return $\bar{s}' = \{(o_i, \tau_i, \mathbf{p}'_i)\}$.
\end{enumerate}
\end{definition}

\begin{remark}[Causal ordering]
\label{rem:ordering}
The ordering in step 2 matters when object responses depend on one another
(e.g., the player must move before the box can detect a push).  In our
implementation, objects read the pre-action context $\text{ctx}(\bar{s})$,
not a partially updated state.  This makes the transition order-independent
for most domains.  For domains where cascading effects require multi-pass
updates (e.g., chain-pushing in Sokoban), the \texttt{respond()} method
returns a \emph{pending action} that is resolved in a second pass.
\end{remark}

\subsection{Relational Preconditions}
\label{sec:relational}

Many game mechanics are \emph{relational}: their effect depends on the spatial
relationship between two or more objects.

\begin{example}[Relational effects]
\label{ex:relational}
\begin{itemize}[nosep]
    \item A box moves only when the player pushes into it (\emph{adjacency} +
          \emph{action direction}).
    \item A switch activates only when an object is on top of it
          (\emph{co-location}).
    \item A key opens a door only when the player is adjacent to the door
          \emph{and} holding the key (\emph{adjacency} + \emph{inventory state}).
\end{itemize}
\end{example}

These effects cannot be discovered by testing objects in isolation.  The
\texttt{respond()} method receives the full relational context, allowing it
to condition on spatial relationships:

\begin{verbatim}
class Box(GameObject):
    def respond(self, action, ctx):
        player = ctx.player()
        if ctx.adjacent(self, opposite(action)) == player:
            push_target = self.position + DELTA[action]
            if ctx.objects_at(push_target) are all non-solid:
                self.position = push_target
\end{verbatim}

The relational predicates available via the context are domain-general
(adjacency, co-location, type queries).  Domain-specific relational
preconditions (e.g., ``same color,'' ``moving toward'') are encoded within
the \texttt{respond()} method by the LLM during synthesis.

\subsection{Connection to STRIPS}
\label{sec:strips_connection}

\begin{remark}[Subsumption of STRIPS]
The object-centric transition subsumes STRIPS~\cite{fikes1971strips}.  Each
$f_{\tau_i}^{\text{respond}}$ encodes the conditional effects of action $a$ on
objects of type $\tau_i$.  The add and delete lists are implicit in the
difference between $\mathbf{p}_i$ and $\mathbf{p}'_i$.  The advantage over
STRIPS is that conditional effects are factored by object type (via
polymorphism) rather than enumerated in a flat list, and that the Python
encoding permits arbitrary computation (spatial reasoning, iteration, complex
conditionals) that PDDL cannot express.
\end{remark}


% ============================================================================
% SECTION 5: CONTINUOUS EPISTEMIC MODEL
% ============================================================================

\section{Continuous Epistemic Model}
\label{sec:epistemic}

WorldCoder's exploration is driven by optimism ($\phi_2$): the model must
admit a plan to reward.  This is effective but coarse.  We augment it with a
continuous epistemic model that tracks \emph{how well} each aspect of the
world model is understood.

\subsection{The Knowledge Matrix}
\label{sec:knowledge_matrix}

\begin{definition}[Observation Context]
\label{def:obs_context}
An \textbf{observation context} $\xi$ for an observed transition
$(s, a, s')$ is a fingerprint of the spatial and relational configuration at
the time of observation:
\begin{equation}
    \xi = (\{\tau_j : o_j \in \text{neighbors}(o_i)\}, \;
           \text{relative positions}, \;
           \text{selection state})
\end{equation}
Two observations have \emph{different} contexts if the set of neighboring
object types or their relative positions differ.
\end{definition}

\begin{definition}[Transition Record]
\label{def:transition_record}
A \textbf{transition record} for object type $\tau$ and action $a$ is:
\begin{equation}
    r = (\tau, a, \Delta\mathbf{p}, \xi),
\end{equation}
where $\Delta\mathbf{p}$ is the observed change in atomic predicates and
$\xi$ is the observation context.
\end{definition}

\begin{definition}[Cell Knowledge]
\label{def:cell_knowledge}
For each $(\tau, a)$ pair, the \textbf{cell knowledge} is:
\begin{equation}
    \calK_{\tau,a} = (n, \; d, \; c, \; \text{cond}),
\end{equation}
where:
\begin{itemize}[nosep]
    \item $n \in \bbN_0$: \textbf{observation count}.
    \item $d \in \bbN_0$: \textbf{context diversity} (number of distinct
          contexts $\xi$ observed).
    \item $c \in [0, 1]$: \textbf{effect consistency} (fraction of
          observations agreeing with the majority effect).
    \item $\text{cond} \in \{0, 1\}$: \textbf{conditional flag}, set when
          the same $(\tau, a)$ pair produces different effects in different
          contexts.
\end{itemize}
\end{definition}

\begin{definition}[Epistemic State]
\label{def:epistemic_state}
The \textbf{epistemic state} is the matrix
$\calE = [\calK_{\tau_i, a_j}]_{i \leq K, j \leq N}$,
where $K = |\calV|$ is the number of object types and $N = |\calA|$ is the
number of actions.
\end{definition}

\subsection{Exploration Priority}
\label{sec:priority}

\begin{definition}[Exploration Priority]
\label{def:priority}
The \textbf{exploration priority} of cell $(\tau, a)$ is:
\begin{equation}
\label{eq:priority}
    \pi(\tau, a) = \frac{w_1}{1 + n_{\tau,a}}
                 + \frac{w_2}{1 + d_{\tau,a}}
                 + w_3 \cdot (1 - c_{\tau,a}),
\end{equation}
where $w_1, w_2, w_3 > 0$ are weights.  The three terms reward:
(1) under-tested cells, (2) low context diversity, and (3) inconsistent
observations (indicating unidentified preconditions).
\end{definition}

\begin{remark}[Comparison to WorldCoder's exploration]
WorldCoder's $\phi_2$ drives exploration by requiring that the model be
optimistic about unvisited regions.  This is a global property of the
model.  Our exploration priority is \emph{local} to each (type, action) cell,
enabling targeted experiments.  The two mechanisms are complementary: $\phi_2$
determines \emph{where to go} (which region of state space), and $\pi(\tau, a)$
determines \emph{what to test} (which type-action interaction).
\end{remark}

\subsection{The Conflation Problem}
\label{sec:conflation}

A binary known/unknown partition has a fundamental failure mode:

\begin{example}[Blocked action recorded as no-effect]
\label{ex:conflation}
The player attempts to move right, but a wall is adjacent.  No visible effect
occurs.  A binary model records ``ACTION\_RIGHT has no effect on the player''
(Known False).  But ACTION\_RIGHT \emph{does} move the player---it was
blocked by context.
\end{example}

The continuous model avoids this: it records the observation \emph{with its
context} $\xi$ (wall was adjacent in direction of movement).  When a later
observation in a different context (open space) shows the player moving, the
cell's consistency $c$ drops below 1, the conditional flag is set, and the
exploration priority increases.  The agent then actively seeks diverse contexts
to identify the precondition.

\subsection{Active Hypothesis Testing}
\label{sec:hypothesis_testing}

When a cell $(\tau, a)$ has $\text{cond} = 1$ (conditional effect detected),
the exploration engine enters a \emph{relational probing} mode:

\begin{enumerate}[nosep]
    \item Examine the observation records for the cell.
    \item Identify the contextual feature that differs between consistent and
          inconsistent observations (e.g., presence/absence of an adjacent
          object of a specific type).
    \item Design an experiment: navigate to a state where the hypothesized
          precondition holds, execute the action, and observe.
    \item If the hypothesis is confirmed, the \texttt{respond()} method for
          $\tau$ should encode the conditional.  If refuted, generate a new
          hypothesis.
\end{enumerate}

This is a form of active learning~\cite{settles2009active} applied to
transition dynamics: the agent selects the most informative state-action
pair, not just the most uncertain one.


% ============================================================================
% SECTION 6: FACTORED SYNTHESIS
% ============================================================================

\section{Factored Program Synthesis}
\label{sec:synthesis}

\subsection{The Factored Model Class}
\label{sec:factored}

\begin{definition}[OOP World Model]
\label{def:oop_world_model}
An \textbf{OOP world model} is:
\begin{equation}
    \hat{W} = (\calH, \; \{f_{\tau_k}^{\text{respond}}\}_{k=1}^K, \;
               f_{\text{perceive}}),
\end{equation}
where $\calH$ is the type hierarchy, $f_{\tau_k}^{\text{respond}}$ is the
response method for type $\tau_k$, and $f_{\text{perceive}}$ is the (fixed)
perception function.  The induced transition function is:
\begin{equation}
    \hat{T}_{\hat{W}}(s, a) = \text{Render}\Bigl(
        \bigl\{f_{\tau_i}^{\text{respond}}(o_i, a, \text{ctx}(\bar{s}))
        \bigr\}_{i=1}^M \Bigr),
    \qquad \bar{s} = f_{\text{perceive}}(s).
\end{equation}
\end{definition}

\begin{definition}[Factored Model Class]
\label{def:factored_class}
Let $\calV = \{\tau_1, \ldots, \tau_K\}$ be the object types.  The
\textbf{factored model class} is:
\begin{equation}
    \calW_{\text{OOP}} = \prod_{k=1}^{K} \calW_{\tau_k},
\end{equation}
where $\calW_{\tau_k}$ is the set of possible implementations of type $\tau_k$'s
response method.  Perception is fixed and excluded from the product.
\end{definition}

Compare with WorldCoder's monolithic model class $\calW_{\text{mono}}$, which
is the set of all possible implementations of the single
\texttt{transition\_function}.

\subsection{Sample Complexity Bound}
\label{sec:bound}

\begin{theorem}[Factored CEGIS Bound]
\label{thm:factored_bound}
Under the OOP factorization, the total number of counterexamples required to
converge to the correct world model is at most:
\begin{equation}
    K_{\bot\!\!\bot}(\calW_{\text{OOP}}) \leq
        \sum_{k=1}^{K} K_{\bot\!\!\bot}(\calW_{\tau_k}).
\end{equation}
The total interaction count is at most:
\begin{equation}
    D \cdot \left(\sum_{k=1}^{K}
        K_{\bot\!\!\bot}(\calW_{\tau_k}) + 1\right),
\end{equation}
where $D$ is the environment diameter.
\end{theorem}

\begin{proof}[Proof sketch]
A counterexample $(s, a, s')$ where $\hat{T}(s, a) \neq s'$ can be attributed
to a specific object type: the type $\tau_k$ such that
$f_{\tau_k}^{\text{respond}}(o_k, a, \text{ctx}) \neq o'_k$ where $o'_k$ is
the true post-transition state of object $o_k$.  Since response methods are
independent, each counterexample eliminates at least one candidate from one
type's hypothesis space.  The total eliminations are bounded by the sum.
\end{proof}

\begin{corollary}[Advantage over monolithic]
\label{cor:advantage}
Since each $\calW_{\tau_k}$ is a strict subset of $\calW_{\text{mono}}$
(it only specifies one type's behavior), we have
$K_{\bot\!\!\bot}(\calW_{\tau_k}) \ll K_{\bot\!\!\bot}(\calW_{\text{mono}})$.
Concretely, if each type has comparable complexity and the monolithic function
has no structural constraints, the monolithic dimensionality can be as large as
$\prod_k |\calW_{\tau_k}|$ in the worst case, versus $\sum_k |\calW_{\tau_k}|$
for the factored model.
\end{corollary}

\subsection{Per-Type CEGIS Loop}
\label{sec:per_type_cegis}

The synthesis loop operates at the per-type level:

\begin{algorithm}[h]
\caption{Per-Type CEGIS Synthesis}
\label{alg:synthesis}
\begin{algorithmic}[1]
\STATE \textbf{Input:} Replay buffer $\calD$, current model $\hat{W}$,
       epistemic state $\calE$
\STATE \textbf{Select target type:} $\tau^* \leftarrow
       \argmax_{\tau \in \calV} \text{synthesis\_priority}(\tau, \calE, \calD)$
\STATE \textbf{Extract:} $\calD_{\tau^*} \leftarrow \{(s, a, s') \in \calD :
       \exists o \in s \text{ of type } \tau^* \text{ with }
       f_{\tau^*}^{\text{respond}}(o, a, \text{ctx}) \neq o'\}$
       \COMMENT{counterexamples for $\tau^*$}
\STATE \textbf{Prepare context:} Include current $f_{\tau^*}^{\text{respond}}$
       code, counterexamples $\calD_{\tau^*}$, and signatures of other types'
       respond methods
\REPEAT
    \STATE Prompt LLM to generate/edit $f_{\tau^*}^{\text{respond}}$
    \STATE Verify against \emph{full} replay buffer $\calD$
    \STATE If new counterexamples: add to prompt as error signal
\UNTIL{all counterexamples resolved \textbf{or} max iterations}
\STATE \textbf{Regression gate:} accept only if
       $\text{Acc}(\hat{W}'; \calD) \geq \text{Acc}(\hat{W}; \calD)$
\end{algorithmic}
\end{algorithm}

The key property: synthesis reasons about one type at a time while verifying
against all transitions.  This focuses the LLM's synthesis budget on the
component that is actually wrong, rather than asking it to reason about all
types simultaneously.

\subsection{Transfer via Inheritance}
\label{sec:transfer}

\begin{definition}[Minimal Subclass Specialization]
\label{def:transfer}
When encountering a new domain, the agent attempts to reuse existing type
classes.  If type $\tau$ from domain A does not match observations in domain B,
the agent derives a subclass $\tau'$ with $\tau'_{\text{parent}} = \tau$ that
overrides only the methods needed to explain the new observations.
\end{definition}

\begin{proposition}[Transfer Sample Complexity]
\label{prop:transfer}
If the correct model for domain B can be obtained from domain A's model by
adding $k$ new subclasses and overriding $r$ response methods total, then the
sample complexity for learning domain B given domain A's model is:
\begin{equation}
    D_{B} \cdot \left(\sum_{i=1}^{k}
        K_{\bot\!\!\bot}(\calW_{\tau'_i}) + r + 1\right),
\end{equation}
compared to $D_{B} \cdot (K_{\bot\!\!\bot}(\calW_{B}) + 1)$ from scratch.
\end{proposition}

\begin{example}[Sokoban $\to$ Sokoban with ice]
\label{ex:transfer}
In standard Sokoban, \texttt{Box.respond()} moves one cell when pushed.  In
Sokoban-with-ice, boxes slide until hitting an obstacle.  The agent derives
\texttt{IceBox} from \texttt{Box}, overriding only the movement logic.
\texttt{Wall}, \texttt{Player}, and \texttt{Goal} are reused without
modification.  Synthesis cost: one class override vs.\ a full transition
function from scratch.
\end{example}


% ============================================================================
% SECTION 7: AGENTIC SYNTHESIS
% ============================================================================

\section{Agentic Synthesis Backend}
\label{sec:agentic}

WorldCoder implements its CEGIS loop via \emph{stateless backprompting}: each
synthesis iteration is an independent LLM API call that receives the full
context (replay buffer, prior code, counterexamples) as a prompt and returns
a complete candidate program as text.  The system then verifies the program
externally and constructs a new prompt for the next iteration.

We replace this with an \emph{agentic synthesis backend}: the LLM operates as
an autonomous coding agent with access to a persistent workspace on disk.

\subsection{Stateless Backprompting vs.\ Agentic Synthesis}
\label{sec:stateless_vs_agentic}

\begin{definition}[Stateless CEGIS Iteration]
\label{def:stateless_cegis}
In WorldCoder, each CEGIS iteration $k$ is a function:
\begin{equation}
    \theta^{(k)} = \text{LLM}(\calD, \theta^{(k-1)}, E^{(k-1)}),
\end{equation}
where $\calD$ is the replay buffer, $\theta^{(k-1)}$ is the prior program, and
$E^{(k-1)}$ is the set of counterexamples from verification.  The LLM call is
\textbf{stateless}: the model has no memory beyond what is provided in the
prompt, and returns a complete program as a single text output.
\end{definition}

\begin{definition}[Agentic CEGIS Iteration]
\label{def:agentic_cegis}
In our system, a single synthesis invocation launches an autonomous agent
with access to a \textbf{persistent workspace} $W$:
\begin{equation}
    W = (\texttt{game\_engine.py}, \; \texttt{test\_runner.py}, \;
         \texttt{context.txt}, \; \texttt{replay\_buffer.pkl}, \;
         \texttt{history.json}).
\end{equation}
The agent executes an \emph{internal} CEGIS loop within a single invocation:
\begin{enumerate}[nosep]
    \item Read \texttt{context.txt} (structured world model description) and
          the current \texttt{game\_engine.py}.
    \item Edit \texttt{game\_engine.py} (incremental modifications, not full
          rewrites).
    \item Execute \texttt{test\_runner.py} via shell (runs the synthesized code
          against the replay buffer).
    \item Read test output (per-transition pass/fail, predicted vs.\ actual
          object positions).
    \item Diagnose failures, edit targeted code, re-run tests.
    \item Repeat until all tests pass or the agent's turn budget is exhausted.
\end{enumerate}
\end{definition}

The key differences are:

\begin{table}[h]
\centering
\caption{Stateless backprompting vs.\ agentic synthesis.}
\label{tab:agentic}
\begin{tabular}{@{}lll@{}}
\toprule
& \textbf{WorldCoder} & \textbf{Ours} \\
\midrule
LLM interaction & Stateless API calls & Persistent interactive session \\
Code editing & Whole-program replacement & Incremental file edits \\
Error diagnosis & Counterexamples in prompt & Agent reads test output, reasons \\
Tool access & None (text $\to$ text) & File system + shell execution \\
Memory & Only what fits in prompt & Full conversation + disk state \\
Strategy & Fixed: synthesize $\to$ verify & Agent chooses: debug, refactor, revert \\
\bottomrule
\end{tabular}
\end{table}

\subsection{Persistent Workspace}
\label{sec:workspace}

The workspace persists across synthesis invocations (not just within one).
When new observations are collected and synthesis is re-triggered, the agent
sees its previous code and can edit incrementally rather than rewriting from
scratch.  This enables:

\begin{itemize}[nosep]
    \item \textbf{Incremental refinement}: new observations refine existing
          code rather than replacing it.
    \item \textbf{History tracking}: \texttt{history.json} records accuracy
          over synthesis iterations, enabling the agent to detect regressions.
    \item \textbf{Structured error signals}: the test runner reports exactly
          which objects predicted incorrectly and by how much (position deltas),
          enabling targeted fixes rather than wholesale rewrites.
\end{itemize}

\subsection{Synthesis Backend Abstraction}
\label{sec:backend_abstraction}

The agentic synthesis is implemented behind an abstract
\texttt{SynthesisBackend} interface: any system that can read a workspace,
edit \texttt{game\_engine.py}, run \texttt{test\_runner.py}, and iterate on
failures can serve as the backend.  The primary implementation uses an
interactive coding agent (Claude Code~\cite{claudecode}) spawned as a
subprocess, but the interface supports any agentic coding system, API-based
reflexion, or even manual human editing.

\begin{remark}[Both are CEGIS]
Both WorldCoder's backprompting and our agentic synthesis implement the CEGIS
pattern (synthesize $\to$ verify $\to$ refine on counterexamples).  The
difference is the \emph{mechanism}: WorldCoder executes CEGIS via prompt
engineering over a stateless API, while our system executes CEGIS via an
autonomous agent with tools.  The agent can perform multi-step reasoning
within a single CEGIS iteration---adding debug prints, testing hypotheses
about failure causes, refactoring code---that would require multiple
stateless iterations in WorldCoder.
\end{remark}


% ============================================================================
% SECTION 8: JOINT LEARNING: HIERARCHY + MODEL
% ============================================================================

\section{Joint Learning: Hierarchy and Model}
\label{sec:joint}

The agent must simultaneously learn \emph{what object types exist} (the
hierarchy $\calH$) and \emph{how they behave} (the response methods).  These
are coupled: identifying types requires observing behavior, and expressing
behavior requires a type vocabulary.

\subsection{Two Types of Counterexamples}
\label{sec:cx_types}

\begin{definition}[Model vs.\ Representation Counterexamples]
\label{def:cx_types}
A counterexample $(s, a, s')$ where $\hat{T}(s, a) \neq s'$ is:
\begin{enumerate}[nosep]
    \item A \textbf{model counterexample} if the type hierarchy is adequate
          (a correct implementation exists within the current factored model
          class $\calW_{\text{OOP}}$).  Fix: re-synthesize the faulty
          response method.
    \item A \textbf{representation counterexample} if no implementation within
          $\calW_{\text{OOP}}$ resolves it.  Fix: grow the hierarchy (add new
          types, new predicates, or subclass specializations).
\end{enumerate}
\end{definition}

\begin{proposition}[Distinguishing counterexample types]
\label{prop:distinguish}
If synthesis fails to resolve a counterexample after $k_{\max}$ attempts,
declare it a representation counterexample and trigger hierarchy refinement.
Model counterexamples are resolved by re-synthesis; representation
counterexamples persist across synthesis attempts.
\end{proposition}

\subsection{The Learning Algorithm}
\label{sec:algorithm}

\begin{algorithm}[h]
\caption{OOP World Model Learning (Full System)}
\label{alg:main}
\begin{algorithmic}[1]
\STATE \textbf{Initialize:} $\calH \leftarrow \{\texttt{GameObject}\}$,
       $\calD \leftarrow \emptyset$, $\calE \leftarrow \text{empty matrix}$
\REPEAT
    \STATE \textbf{// Exploration}
    \STATE Select $(\tau^*, a^*)$ via exploration priority
           (Eq.~\ref{eq:priority})
    \STATE Navigate to an instance of $\tau^*$; execute $a^*$; observe
           $(s_t, a^*, s_{t+1})$
    \STATE Update $\calD$, $\calE$
    \STATE \textbf{// Verification}
    \IF{$\hat{T}(s_t, a^*) \neq s_{t+1}$}
        \STATE Attribute error to object type $\tau_k$
        \STATE \textbf{// Inner loop: per-type CEGIS (Algorithm~\ref{alg:synthesis})}
        \STATE Re-synthesize $f_{\tau_k}^{\text{respond}}$
        \IF{counterexample persists after $k_{\max}$ attempts}
            \STATE \textbf{// Hierarchy refinement}
            \STATE Propose new type or predicate; grow $\calH$
        \ENDIF
    \ENDIF
\UNTIL{convergence criterion met}
\end{algorithmic}
\end{algorithm}

\subsection{Sample Complexity of Joint Learning}
\label{sec:joint_bound}

\begin{theorem}[Joint Learning Bound]
\label{thm:joint}
Let $\calH^*$ be the true adequate type hierarchy and $\hat{W}^*$ the true
world model.  Suppose the synthesis oracle succeeds within $k_{\max}$ attempts
for model counterexamples and the total number of distinct predicates/types
ever considered is $|\Psi_{\text{total}}|$.  Then the total interactions are
at most:
\begin{equation}
    D \cdot \left(
        \underbrace{|\Psi_{\text{total}}|}_{\text{hierarchy cost}}
        + \underbrace{\sum_{\tau \in \calH^*}
        K_{\bot\!\!\bot}(\calW_\tau)}_{\text{model cost}}
        + 1
    \right).
\end{equation}
\end{theorem}

\begin{proof}[Proof sketch]
The outer loop (hierarchy refinement) iterates at most $|\Psi_{\text{total}}|$
times.  Each inner loop iteration requires at most
$D \cdot (\sum_k K_{\bot\!\!\bot}(\calW_{\tau_k}) + 1)$ interactions by
Theorem~\ref{thm:factored_bound}.  Counterexamples accumulate across
iterations, so they need not be re-discovered.
\end{proof}


% ============================================================================
% SECTION 8: EXPLORATION
% ============================================================================

\section{Active Exploration}
\label{sec:exploration}

\subsection{Phased Exploration}
\label{sec:phases}

Exploration proceeds through five phases:

\begin{enumerate}[nosep]
    \item \textbf{DISCOVER}: Identify object types from initial observations.
    \item \textbf{IDENTIFY\_AGENT}: Determine which object the agent controls
          by testing directional actions.
    \item \textbf{SURVEY}: Execute each action on each reachable object type
          once, populating the epistemic matrix.
    \item \textbf{TARGETED}: Priority-driven testing (Eq.~\ref{eq:priority}):
          select the highest-priority $(\tau, a)$ cell, navigate to an instance
          of $\tau$, execute $a$, observe.
    \item \textbf{RELATIONAL}: When conditional effects are detected, probe
          interactions between object pairs by systematically varying spatial
          context.
\end{enumerate}

Phase transitions are automatic, based on the epistemic state.

\subsection{Integration with WorldCoder's Optimism}
\label{sec:optimism}

Our exploration priority (Eq.~\ref{eq:priority}) complements WorldCoder's
optimism constraint ($\phi_2$):

\begin{itemize}[nosep]
    \item $\phi_2$ determines the \textbf{high-level exploration goal}: plan
          toward unvisited state regions using the current model.
    \item $\pi(\tau, a)$ determines the \textbf{low-level experimental target}:
          which (type, action) interaction to test next.
\end{itemize}

When $\phi_2$ generates a plan that traverses unknown territory, the agent
executes the plan.  At each step, it updates the epistemic state.  When the
plan completes or reaches an impasse, the agent falls back to
priority-driven targeted exploration.

\subsection{Setup Sequences and Navigation}
\label{sec:setup}

To reach informative states, the agent constructs \textbf{setup sequences}
of verified actions (actions whose response methods have high confidence in
the epistemic state).  For grid domains, this includes a
``move to sprite'' macro primitive: BFS pathfinding to a target object's
position, composed from directional actions.

Only actions with consistency $c_{\tau,a} > \theta_{\text{safe}}$ are used
in setup sequences, providing an analogous safety guarantee to the safe
action models of Juba et al.~\cite{juba2021safe}.


% ============================================================================
% SECTION 9: EVALUATION
% ============================================================================

\section{Evaluation}
\label{sec:evaluation}

\subsection{Benchmarks}
\label{sec:benchmarks}

\begin{table}[h]
\centering
\caption{Evaluation benchmarks.}
\label{tab:benchmarks}
\begin{tabular}{@{}llllp{4cm}@{}}
\toprule
\textbf{Benchmark} & \textbf{State} & \textbf{Actions} & \textbf{Types} &
\textbf{Key Challenge} \\
\midrule
Sokoban & Symbolic & 4 & 4 & Push chains, deadlocks \\
MiniGrid & Symbolic & 7 & $\sim$6 & Transfer across variants \\
AlfWorld & Text & NL & $\sim$20 & Large type space, compositional \\
ARC-AGI-3 & 64$\times$64 RGB & 5--6 & varies & Raw perception, anonymous
actions, diverse mechanics \\
\bottomrule
\end{tabular}
\end{table}

\subsection{Metrics}
\label{sec:metrics}

\textbf{Primary:} Transition accuracy---percentage of observed transitions
correctly predicted by $\hat{T}_{\hat{W}}$:
\begin{equation}
    \text{Acc}(\hat{T}; \calD) = \frac{1}{|\calD|}
        \sum_{(s,a,s') \in \calD} \ind[\hat{T}(s, a) = s'].
\end{equation}

\textbf{Secondary:}
\begin{itemize}[nosep]
    \item \textbf{Sample efficiency}: interactions to reach $X$\% accuracy.
    \item \textbf{Synthesis efficiency}: LLM tokens consumed.
    \item \textbf{Transfer score}: accuracy on domain B using domain A's model.
\end{itemize}

\subsection{ARC-AGI-3 Perception Adapter}
\label{sec:arc_perception}

For Sokoban, MiniGrid, and AlfWorld, perception is trivial: the environment
provides structured object state directly, as in WorldCoder.  ARC-AGI-3 is
the only benchmark requiring a perception layer, since the environment returns
raw 64$\times$64 RGB frames.

We implement a deterministic perception pipeline as a compatibility layer:
\begin{enumerate}[nosep]
    \item \textbf{Fragment detection}: color-based connected-component labeling
          with 8-connectivity on each non-background color channel
          (\texttt{scipy.ndimage.label}).
    \item \textbf{Sprite registry}: maintains a library of discovered sprite
          types, matched by color and shape.  Supports merge (when co-movement
          evidence indicates multi-color composites) and split.
    \item \textbf{Co-movement tracking}: discovers multi-color sprites by
          observing which fragments consistently displace by the same vector
          across frames.
    \item \textbf{Object tracking}: assigns persistent track IDs across frames
          via a three-pass matching algorithm (color match $\to$ spatial overlap
          $\to$ nearest-neighbor centroid).
\end{enumerate}

This pipeline is deterministic, parameter-free, and not part of the synthesis
target.  It is functionally equivalent to the structured state that Sokoban,
MiniGrid, and AlfWorld provide natively---it simply bridges the gap for
pixel-observation environments.

\subsection{Baselines and Ablations}
\label{sec:baselines}

\begin{enumerate}[nosep]
    \item \textbf{WorldCoder}~\cite{tang2024worldcoder}: monolithic CEGIS
          with stateless backprompting.
          Direct comparison on all four benchmarks.
    \item \textbf{Ours (monolithic ablation)}: our system with OOP disabled
          (single \texttt{transition\_function}).  Isolates the value of
          factored synthesis.
    \item \textbf{Ours (binary epistemic ablation)}: continuous confidence
          replaced by Known-True/Known-False/Unknown.  Isolates the value
          of the continuous epistemic model.
    \item \textbf{Ours (random exploration ablation)}: priority-driven
          exploration replaced by uniform random action selection.  Isolates
          the value of smart exploration.
    \item \textbf{Ours (stateless synthesis ablation)}: agentic backend
          replaced by stateless API reflexion (single-shot LLM calls with
          backprompting).  Isolates the value of agentic synthesis.
\end{enumerate}

\subsection{Hypotheses}
\label{sec:hypotheses}

\begin{enumerate}[nosep]
    \item \textbf{H1 (Factored synthesis):} On domains with $\geq$3 object
          types, OOP factoring reaches equivalent accuracy with fewer LLM
          tokens than monolithic synthesis.
    \item \textbf{H2 (Exploration efficiency):} Continuous epistemic
          exploration reaches equivalent accuracy with fewer environment
          interactions than WorldCoder's optimism-only exploration.
    \item \textbf{H3 (Transfer):} Learned object classes transfer to novel
          domain variants, reducing synthesis cost.
    \item \textbf{H4 (Agentic synthesis):} The agentic backend reaches higher
          accuracy than stateless backprompting given the same observations
          and synthesis budget.
    \item \textbf{H5 (Pixel observations):} The system achieves $>$80\%
          transition accuracy on ARC-AGI-3, extending beyond symbolic-state
          domains.
\end{enumerate}

\subsection{Preliminary Results}
\label{sec:results}

\begin{table}[h]
\centering
\caption{Preliminary results on ARC-AGI-3 game ls20.  The three rows
isolate the effect of perception strategy and synthesis backend.}
\label{tab:results}
\begin{tabular}{@{}lcc@{}}
\toprule
\textbf{Configuration} & \textbf{Transition Accuracy} & \textbf{Iterations} \\
\midrule
LLM perceive + transition (monolithic) & 0\% & all spent on perception bugs \\
Deterministic perception + stateless API reflexion & 88\% & 4 \\
Deterministic perception + agentic synthesis & 100\% (8/8) & 1 \\
\bottomrule
\end{tabular}
\end{table}

The progression 0\% $\to$ 88\% $\to$ 100\% validates two architectural
decisions.  First, fixing perception as a deterministic preprocessor
(removing it from the synthesis target) allows the LLM to focus entirely on
transition logic.  Second, the agentic backend (H4) reaches 100\% in a single
invocation where stateless API reflexion plateaus at 88\% after 4 iterations,
demonstrating that persistent workspace access and multi-step debugging
provide a concrete advantage over backprompting.


% ============================================================================
% SECTION 10: RELATED WORK
% ============================================================================

\section{Related Work}
\label{sec:related}

\textbf{World model synthesis.}
WorldCoder~\cite{tang2024worldcoder} synthesizes monolithic Python world models
via CEGIS with an optimism constraint.  PoE-World~\cite{poeworld2025} extends
this with a product of programmatic experts.  Our OOP factorization is
complementary: we achieve compositionality through class hierarchy rather than
expert mixture, and additionally support transfer via inheritance.

\textbf{Action model learning.}
Classical approaches learn STRIPS-style schemas from observed
transitions~\cite{cresswell2011acquiring, arora2018review, juba2021safe}.
OLAM~\cite{lamanna2021olam} addresses active exploration for online action
model learning.  Our system learns equivalent information but encodes it as
executable Python with object-centric factoring.

\textbf{Object-oriented MDPs.}
Diuk et al.~\cite{diuk2008oomdp} formalize OO-MDPs with object classes and
relations.  We extend this to the \emph{synthesis} setting: not only is the
state object-oriented, but the transition function itself is factored by
object type.

\textbf{Program synthesis.}
DreamCoder~\cite{ellis2021dreamcoder} bootstraps program synthesis with
library learning.  Our OOP type hierarchy serves an analogous role: learned
classes form a reusable library.
Reflexion~\cite{shinn2024reflexion} uses verbal self-reflection for
LLM agent learning; our synthesis loop uses the same generate-test-reflect
pattern with structured per-object error signals.

\textbf{Object-centric representation learning.}
Slot Attention~\cite{locatello2020object},
C-SWM~\cite{kipf2020contrastive}, and OP3~\cite{veerapaneni2020entity}
learn object representations from pixels via neural networks.  For
grid environments, we use deterministic connected-component detection,
gaining determinism and zero-shot transfer at the cost of generality.

\textbf{Exploration in model-based RL.}
R-MAX~\cite{brafman2002rmax} maintains known/unknown partitions.
VIME~\cite{houthooft2016vime} maximizes information gain.
Go-Explore~\cite{ecoffet2021first} archives interesting states.
Our continuous epistemic model draws on all three: it tracks knowledge at
the (type, action) level, prioritizes information gain, and uses setup
sequences to return to informative states.


% ============================================================================
% SECTION 11: CONCLUSION
% ============================================================================

\section{Conclusion}
\label{sec:conclusion}

We have presented a system that extends WorldCoder's LLM-synthesized world
models with four contributions: OOP-factored synthesis that reduces the CEGIS
search space from monolithic to additive, a continuous epistemic model that
drives targeted exploration, object-centric transition functions with
relational context, and an agentic synthesis backend that replaces stateless
backprompting with persistent, tool-equipped, multi-step reasoning.  The OOP
structure enables cross-domain transfer via class inheritance---a capability
absent from WorldCoder and PoE-World.

Preliminary results on ARC-AGI-3 validate the architecture: the agentic
backend achieves 100\% transition accuracy where stateless reflexion plateaus
at 88\%, and OOP factoring enables per-type debugging that monolithic
synthesis cannot provide.  Full evaluation on Sokoban, MiniGrid, AlfWorld,
and ARC-AGI-3 will test the hypotheses that factored synthesis is more
token-efficient (H1), continuous epistemic exploration is more
sample-efficient (H2), learned classes transfer (H3), agentic synthesis
outperforms backprompting (H4), and the system generalizes to pixel
observations (H5).

\textbf{Limitations.}  The current system assumes deterministic, fully
observable environments.  Partial observability (MiniGrid's limited view),
stochastic transitions, and non-grid state spaces remain open challenges.
The relational context provides a fixed set of spatial queries; discovering
novel relational predicates (e.g., ``same color,'' ``moving toward'') during
synthesis is not yet automated.  The agentic synthesis backend is more
expensive per invocation than stateless API calls; the tradeoff between
per-iteration cost and total iterations to convergence needs empirical
characterization.


% ============================================================================
% BIBLIOGRAPHY
% ============================================================================

\bibliographystyle{plain}

\begin{thebibliography}{30}

\bibitem{arcagi3}
{ARC Prize Foundation}.
\newblock ARC-AGI-3: A Benchmark for General Intelligence.
\newblock \url{https://arcprize.org}, 2025.

\bibitem{arora2018review}
A.~Arora, H.~Fiorino, D.~Pellier, M.~M{\'e}tivier, and S.~Pesty.
\newblock A Review of Learning Planning Action Models.
\newblock \emph{The Knowledge Engineering Review}, 33:e20, 2018.

\bibitem{brafman2002rmax}
R.~I.~Brafman and M.~Tennenholtz.
\newblock R-MAX, A General Polynomial Time Algorithm for Near-Optimal
  Reinforcement Learning.
\newblock \emph{JMLR}, 3:213--231, 2002.

\bibitem{claudecode}
{Anthropic}.
\newblock Claude Code: An Agentic Coding Tool.
\newblock \url{https://docs.anthropic.com/en/docs/claude-code}, 2025.

\bibitem{cresswell2011acquiring}
S.~N.~Cresswell, T.~L.~McCluskey, and M.~M.~West.
\newblock Acquiring Planning Domain Models Using LOCM.
\newblock \emph{The Knowledge Engineering Review}, 28(2):195--213, 2013.

\bibitem{diuk2008oomdp}
C.~Diuk, A.~Cohen, and M.~L.~Littman.
\newblock An Object-Oriented Representation for Efficient Reinforcement Learning.
\newblock In \emph{ICML}, 2008.

\bibitem{ecoffet2021first}
A.~Ecoffet, J.~Huizinga, J.~Lehman, K.~O.~Stanley, and J.~Clune.
\newblock First Return, Then Explore.
\newblock \emph{Nature}, 590:580--586, 2021.

\bibitem{ellis2021dreamcoder}
K.~Ellis, C.~Wong, M.~Nye, et~al.
\newblock DreamCoder: Bootstrapping Inductive Program Synthesis with Wake-Sleep
  Library Learning.
\newblock In \emph{PLDI}, 2021.

\bibitem{fikes1971strips}
R.~E.~Fikes and N.~J.~Nilsson.
\newblock STRIPS: A New Approach to the Application of Theorem Proving to
  Problem Solving.
\newblock \emph{Artificial Intelligence}, 2(3--4):189--208, 1971.

\bibitem{fox2003pddl2}
M.~Fox and D.~Long.
\newblock PDDL2.1: An Extension to PDDL for Expressing Temporal Planning
  Domains.
\newblock \emph{JAIR}, 20:61--124, 2003.

\bibitem{ha2018world}
D.~Ha and J.~Schmidhuber.
\newblock World Models.
\newblock \emph{arXiv:1803.10122}, 2018.

\bibitem{hafner2020dreamer}
D.~Hafner, T.~Lillicrap, J.~Ba, and M.~Norouzi.
\newblock Dream to Control: Learning Behaviors by Latent Imagination.
\newblock In \emph{ICLR}, 2020.

\bibitem{houthooft2016vime}
R.~Houthooft, X.~Chen, Y.~Duan, J.~Schulman, F.~De~Turck, and P.~Abbeel.
\newblock VIME: Variational Information Maximizing Exploration.
\newblock In \emph{NeurIPS}, 2016.

\bibitem{juba2021safe}
B.~Juba and R.~S.~Stern.
\newblock Safe Action Model Learning.
\newblock \emph{arXiv:2107.09336}, 2021.

\bibitem{kipf2020contrastive}
T.~Kipf, E.~van~der~Pol, and M.~Welling.
\newblock Contrastive Learning of Structured World Models.
\newblock In \emph{ICLR}, 2020.

\bibitem{lamanna2021olam}
L.~Lamanna, A.~Saetti, L.~Serafini, A.~Gerevini, and P.~Traverso.
\newblock Online Learning of Action Models for PDDL Planning.
\newblock In \emph{IJCAI}, 2021.

\bibitem{locatello2020object}
F.~Locatello, D.~Weissenborn, T.~Unterthiner, et~al.
\newblock Object-Centric Learning with Slot Attention.
\newblock In \emph{NeurIPS}, 2020.

\bibitem{poeworld2025}
K.~Ellis et~al.
\newblock PoE-World: Compositional World Models via Product of Programmatic
  Experts.
\newblock \emph{arXiv:2505.10819}, 2025.

\bibitem{schrittwieser2020muzero}
J.~Schrittwieser, I.~Antonoglou, T.~Hubert, et~al.
\newblock Mastering Atari, Go, Chess and Shogi by Planning with a Learned Model.
\newblock \emph{Nature}, 588:604--609, 2020.

\bibitem{settles2009active}
B.~Settles.
\newblock Active Learning Literature Survey.
\newblock Technical Report 1648, University of Wisconsin--Madison, 2009.

\bibitem{shinn2024reflexion}
N.~Shinn, F.~Cassano, A.~Gopinath, K.~Narasimhan, and S.~Yao.
\newblock Reflexion: Language Agents with Verbal Reinforcement Learning.
\newblock In \emph{NeurIPS}, 2023.

\bibitem{tang2024worldcoder}
H.~Tang, D.~Key, and K.~Ellis.
\newblock WorldCoder, a Model-Based LLM Agent: Building World Models by Writing
  Code and Interacting with the Environment.
\newblock In \emph{NeurIPS}, 2024.

\bibitem{veerapaneni2020entity}
R.~Veerapaneni, J.~D.~Co-Reyes, M.~Chang, M.~Janner, C.~Finn, J.~Wu,
  J.~Tenenbaum, and S.~Levine.
\newblock Entity Abstraction in Visual Model-Based Reinforcement Learning.
\newblock In \emph{CoRL}, 2020.

\end{thebibliography}

\end{document}
